\section{Introduction}

A central problem in four-dimensional quantum topology is to construct a
tractable field-theoretic invariant that detects exotic smooth structure.
For premodular fusion categories, the Crane--Yetter theory
reduces to classical invariants and does not detect exotic smooth structure
\cite{Reutter2023}.  On the other hand, gauge-theoretic invariants have been
successful in detecting smooth structure, and second homology is part of
their input: Donaldson invariants admit insertions from $H_2$
\cite{Donaldson1990}, while Seiberg--Witten invariants are organized by
$\operatorname{Spin}^c$ structures and their classes in second cohomology
\cite{Witten1994}.  A natural idea is therefore to enrich the Crane--Yetter
theory by second-homology data and ask whether the resulting invariant can
detect exotic smoothness.

Yetter proposed precisely such a refinement in a sketch note in 1993.  For a
premodular fusion category $\cC$,\footnote{Yetter's terminology was an
``artinian semisimple tortile category.''  A tortile category is what is now
called a ribbon category.  With the finite set of simple isomorphism classes,
finite semisimplicity, rigidity, and simple tensor unit assumed here, this is
precisely a ribbon fusion category, equivalently a premodular fusion category.}
his sketch constructed state-sum invariants of a $D$-manifold $M$
valued in the complex-valued functions on $H_{D-2}(M;\operatorname{Gr}(\cC))$,
where $\operatorname{Gr}(\cC)$ is the abelian group of monoidal natural
automorphisms of the identity functor \cite{Yetter1993}. In dimension four this gives a
homologically twisted Crane--Yetter state sum. Yetter's note does not provide
a Kirby-diagram handle formula, nor does it determine whether the refinement
detects exotic smooth structure. We address these two questions for the
four-dimensional Crane--Yetter theory; we do not consider the twisted
Turaev--Viro theory. Throughout, $\cC$ is a premodular fusion category over
$\mathbb C$---in particular, it is finite and semisimple---equipped with a
monoidal natural automorphism
\[
  \eta\colon \operatorname{Id}_{\cC}\Longrightarrow
  \operatorname{Id}_{\cC}.
\]
For a closed oriented smooth four-manifold $X$ and $\alpha\in H_2(X;\mathbb Z)$, we
define a scalar $\CY_{\cC,\eta}(X,\alpha)$ from a Kirby diagram presenting
$(X,\alpha)$. Our first main result is that this scalar is an invariant of the
oriented diffeomorphism class of the marked pair $(X,\alpha)$. Thus the
construction gives a genuine premodular, homologically twisted four-manifold
invariant rather than only a formal modification of the state sum.  Motivated
by Tham's handle description of Crane--Yetter theory
\cite[Definition~5.11]{Tham2021}, we expect this invariant
to agree with the four-dimensional state sum sketched by Yetter; that
identification is not needed for, and is not claimed as part of, our results.

The motivating question nevertheless has a negative answer in the simply
connected case.  We prove that $\CY_{\cC,\eta}(X,\alpha)$ depends only on the
marked integral intersection lattice
\[
  \bigl(H_2(X;\mathbb Z),Q_X,\alpha\bigr),
\]
namely the integral second-homology group, its intersection pairing, and the
distinguished class $\alpha$, together with categorical constants.  For an odd
intersection form, the normalized value is a product of shifted positive and
negative Gauss sums evaluated on the diagonal coordinates of $\alpha$
(Proposition~\ref{prop:odd-formula}).  For an even intersection form, it is an
explicit stabilized quotient of the values on $K3$, reversed $K3$, and
hyperbolic blocks (Proposition~\ref{prop:even-formula}).  If $\cC$ is modular,
these two descriptions reduce to a single closed formula
(Proposition~\ref{prop:modular-formula}).  On setting the character twist to
be trivial, the odd and even formulas specialize respectively to the non-spin
and spin formulas of Reutter's Corollary~3.9, equations (5) and (6), for the
Crane--Yetter theory \cite{Reutter2023}.
Consequently this most direct semisimple homological refinement of the
Yetter--Crane--Yetter construction does not furnish an exotic-smoothness
witness for simply connected closed four-manifolds
(Corollary~\ref{cor:no-simply-connected-exotic-detection}).

The restriction to simply connected manifolds enters through the stable
classification and realization results used to pass from a marked
intersection lattice to standard smooth blocks.  Wall's stable theorem and
the one-stable realization theorem of Ruberman--Strle are stated in the simply
connected setting \cite{Wall1964,RubermanStrle2022}.  If analogous marked
stable realization and decomposition theorems were established for a
non-simply-connected class of four-manifolds, the argument of this paper would
extend to that class.  In their absence, our result makes no claim that this
homological refinement fails to detect exotic smooth structures outside the
simply connected setting.  However, we believe that the non-simply-connected
case is the natural remaining setting in which this homological refinement may
still detect exotic smooth structure; we state this expectation precisely in
Conjecture~\ref{conj:non-simply-connected-detection}.

The construction and the no-go calculation share the same concrete starting
point: a handle formula that makes the homological twist visible.  It is
inspired by the turned-over handle description in Tham's Definition~5.11
\cite{Tham2021}.  Let $D$ be a Kirby diagram with one $0$-handle and one
$4$-handle, with dotted components representing the $1$-handles, ordinary
framed components representing the $2$-handles, and with the boundary after
the $2$-handles closed by $h_3(D)$ $3$-handles.  For a two-handle representative
$a=(a_k)$ of $\alpha$, it has the form
\[
  \CY_{\cC,\eta}(D,a)
  = \frac{\operatorname{Dim}(\cC)}
          {\operatorname{Dim}(\cC)^{h_3(D)}}
    \sum_{\lambda}
      \left(\prod_k d_{\lambda_k}\right)
      \left(\prod_k \eta_{\lambda_k}^{a_k}\right)
      \bigl\langle G_{D,\lambda}\bigr\rangle_{\cC}.
\]
Here $\lambda$ ranges over simple labels of the ordinary framed
$2$-handle components, while every dotted $1$-handle is evaluated by the
normalized projector onto the tensor-unit channel.  The closed ribbon graph
$G_{D,\lambda}$ is the resulting role-aware graphical evaluation.  When
$\alpha=0$, the character factor is trivial and the formula recovers the
untwisted Crane--Yetter handle evaluation of
\cite[Definition~5.11]{Tham2021}.  Proposition~\ref{prop:diffeomorphism-invariance}
states the resulting marked diffeomorphism invariant.  Its proof is direct:
we verify invariance under each of the eight elementary Kirby moves.

As in Reutter's calculation of semisimple four-dimensional field theories,
stabilization turns this invariant into an effective computational tool
\cite[Theorem~A and Corollary~3.9]{Reutter2023}.  Put
$H=S^2\times S^2$, write
$\widehat{\CY}=\operatorname{Dim}(\cC)^{-1}\CY$, and set
$h=\widehat{\CY}_{\cC,\eta}(H,0)$.  Wall's theorem says that simply connected
four-manifolds with isometric intersection forms become diffeomorphic after
taking connected sums with copies of $H$ \cite[Theorems~2--3]{Wall1964}.
Because our invariant is marked by a class in $H_2$, we also use
Ruberman--Strle's one-stable realization theorem to arrange that the stable
diffeomorphism carries the chosen class to the chosen class
\cite[Theorem~2.1]{RubermanStrle2022}.  The connected-sum formula and the
direct calculation $h\neq0$ then imply that a marked stable diffeomorphism
\[
  (X\mathbin{\#}rH,\alpha\oplus0)
  \cong (M\mathbin{\#}sH,\beta\oplus0)
\]
determines the unstabilized value by
\[
  \widehat{\CY}_{\cC,\eta}(X,\alpha)
  =\widehat{\CY}_{\cC,\eta}(M,\beta)h^{s-r}.
\]
Thus stabilization is useful here not merely because it identifies more
manifolds, but because its nonzero contribution can be divided out.

This is especially effective for simply connected four-manifolds.  Their
oriented homeomorphism types are controlled by their unimodular intersection
forms, together with the Kirby--Siebenmann invariant, which vanishes in the
smooth case \cite{Freedman1982}.  The odd indefinite forms are diagonal
\cite[Theorem~2.3.1]{PetersSterk2024}, while Donaldson's theorem gives the same
conclusion for smooth definite forms \cite[Theorem~A]{Donaldson1983}; their
standard models are connected sums of $\mathbb{CP}^2$ and
$\overline{\mathbb{CP}}^{\,2}$.  Even forms become, after hyperbolic
stabilization, connected sums of $K3$, $\overline{K3}$, and $H$ blocks, as in
the proof of \cite[Corollary~3.9]{Reutter2023}.  Consequently the preceding
stable calculation reduces the value on an arbitrary simply connected marked
manifold to values on these few standard pieces.  Connected-sum
multiplicativity and stabilization then give the odd and even closed formulas
in Propositions~\ref{prop:odd-formula} and
\ref{prop:even-formula}.  If $\cC$ is modular, the character--object
correspondence and the Gauss-sum identities of
\cite[Lemmas~8.15.3 and~8.22.9, Corollary~8.15.5, and
Proposition~8.15.4]{EGNO2015} reduce both formulas further to
Proposition~\ref{prop:modular-formula}.
